\documentclass[11pt]{article}
\usepackage[margin=1in]{geometry}
\usepackage{booktabs}
\usepackage{hyperref}
\usepackage[T1]{fontenc}

\title{Cluster B R6 Robustness Consolidation}
\author{Six Birds Theory construction track}
\date{2026-06-05}

\begin{document}
\maketitle

\section{Grade}

\textbf{Grade: finite-carrier robustness consolidation.}
This section consolidates verified Cluster B Steps 6--8. The R6 stress battery was run on the declared finite toy carrier and its finite toy extensions. The structural claims are robust under that declared battery; the modeled forms remain bounded toy forms. This is not a frame-transfer certificate. Richer continuous, Lorentzian-field, gauge, diffeomorphism, and real-RG carriers remain the standing gate. This section supplies no new physics, no \(\Lambda\) value, no physical high-curvature mechanism, and no frame-transfer certificate.

Sources: \texttt{steps/step6\_shared\_frame\_robustness\_artifacts/results\_summary.md}, \texttt{steps/step7\_e042\_earned\_continuation\_artifacts/results\_summary.md}, \texttt{steps/step8\_e021\_rg\_measure\_robustness\_artifacts/results\_summary.md}.

\section{Shared-Frame Representation-Independence}

\textbf{Claim R6.1. Grade: finite-toy-diagnostic / representation-independence.}
Step 6 tests \(d4_{\mathrm{subplanck}}\) and \(d5_{\mathrm{vacuum}}\) under five general-linear readouts and three nonlinear readouts of GR's smooth source \((d0,d2,d3)\). The minimum obstruction counts remain positive:
\[
\min \Delta(d4)=3,\qquad \min \Delta(d5)=1.
\]
Genuine endpoint-internal enrichment to \((d0,d1,d2,d3)\), including polynomial and mixed-linear maps, still leaves obstruction counts \(3\) and \(1\). The function-of-source reason is direct: states identical in the source stay identical under any readout of that source, so branch-pair distinctions outside the source cannot be recovered by reparametrizing it.

Sources: \texttt{steps/step6\_shared\_frame\_robustness\_artifacts/readout\_battery\_step6.csv}, \texttt{steps/step6\_shared\_frame\_robustness\_artifacts/enrichment\_step6.csv}.

\textbf{Claim R6.2. Grade: finite-toy-diagnostic / anti-smuggling control.}
Step 6 also tests adversarial endpoint enrichment. Adding \(d4\) or \(d5\) as endpoint content makes the corresponding target definable with obstruction \(0\), but those rows are flagged as \texttt{uses\_layer\_content=True}. Positive-detection controls have teeth: \(d3\) factors with obstruction \(0\), and \(d0\cdot d2\) factors through polynomial GR \(\Sigma_f\) with obstruction \(0\).

Sources: \texttt{steps/step6\_shared\_frame\_robustness\_artifacts/enrichment\_step6.csv}, \texttt{steps/step6\_shared\_frame\_robustness\_artifacts/controls\_step6.csv}.

\section{E042 Earned Continuation}

\textbf{Claim R6.3. Grade: finite-toy-diagnostic / earned rule-computed continuation.}
Step 7 converts the Step 5 R1 weakness. The post-locus continuation is no longer assigned by \(\mathrm{L\_defined=True}\); it is generated by the finite recurrence
\[
x_{\mathrm{next}}=x-\mathrm{step}\cdot K/(K+K_{\mathrm{Planck}}).
\]
The final post-locus \(U_{\mathrm{good}}\) output is \(-0.7821121792445374\). GR's raw-curvature evolution terminates at or before the locus, and \(U_{\mathrm{bad}}\), driven by raw divergent curvature, fails to continue. The bad-rule control shows the finite continuation is not automatic.

Sources: \texttt{steps/step7\_e042\_earned\_continuation\_artifacts/earned\_continuation\_step7.csv}, \texttt{steps/step7\_e042\_earned\_continuation\_artifacts/controls\_step7.csv}.

\textbf{Claim R6.4. Grade: finite-toy-diagnostic / Lorentzian toy stress.}
Step 7 also replaces the earlier simple curvature shape with a Schwarzschild-like toy curvature \(K(r)=48M^2/r^6\). The pre-locus finite value reaches \(3298534883328.0\), the locus row is \(K=\infty\), and \(R_L\) stabilizes to \(0.75\) with final listed value \(0.7499999997671694\). The metric and recurrence are toys, not a diffeomorphism-invariant constraint algebra or real GR dynamics.

Sources: \texttt{steps/step7\_e042\_earned\_continuation\_artifacts/lorentzian\_curvature\_step7.csv}, \texttt{steps/step7\_e042\_earned\_continuation\_artifacts/nonclaim\_boundary.md}.

\section{E021 Selection Relocation}

\textbf{Claim R6.5. Grade: finite-toy-diagnostic / relocation-finding.}
Step 8 replaces the hand-built Step 3 selection pair with a structured finite toy: a moduli grid, deterministic beta-function RG flow, and a measure/relaxation score selecting \(\phi^\ast\). The realized \(\Lambda\) remains non-descending from the bare UV:
\[
\Delta(\Lambda\mid \mathrm{bare\ UV})=4.
\]
When the selected modulus is included in \(\Sigma_f\), \(\Lambda\) descends with obstruction \(0\), but \(\phi^\ast\) remains non-descending from the bare UV with obstruction \(4\). Thus the non-descendingness relocates to the selection rather than disappearing.

Sources: \texttt{steps/step8\_e021\_rg\_measure\_robustness\_artifacts/nonfactorization\_step8.csv}, \texttt{steps/step8\_e021\_rg\_measure\_robustness\_artifacts/rg\_measure\_vacua\_step8.csv}.

\textbf{Claim R6.6. Grade: finite-toy-diagnostic / positive controls.}
Step 8's controls show the test can detect UV derivability. The RG-derived running \(g_{\mathrm{IR}}\) descends from the bare UV with obstruction \(0\), and the no-selection \(\Lambda\)-like control also descends with obstruction \(0\). The E021 typing survives on the toy, located at the selection predicate.

Sources: \texttt{steps/step8\_e021\_rg\_measure\_robustness\_artifacts/controls\_step8.csv}, \texttt{steps/step8\_e021\_rg\_measure\_robustness\_artifacts/nonfactorization\_step8.csv}.

\section{What Survived and What Stays Bounded}

\textbf{Claim R6.7. Grade: finite-carrier robustness synthesis.}
Survived under the declared battery: \(d4/d5\) non-descending is representation-independent under the tested readouts and genuine enrichments; E042 continuation is rule-computed with a bad-rule control; E021 non-descendingness relocates to the selection under adversarial enrichment. Bounded: the \(R_L\) saturation shape, the Lorentzian metric toy, the finite recurrence, and the RG/measure/moduli toy. The R6 battery did not break the structural claims; it bounded the modeled forms.

Sources: \texttt{steps/step6\_shared\_frame\_robustness\_artifacts/results\_summary.md}, \texttt{steps/step7\_e042\_earned\_continuation\_artifacts/results\_summary.md}, \texttt{steps/step8\_e021\_rg\_measure\_robustness\_artifacts/results\_summary.md}.

\section{Relation to the v1 Review}

\textbf{Claim R6.8. Grade: organizational / review-response.}
R1 is converted by Step 7: continuation is earned by a recurrence rather than assigned. R2 was narrowed in Step 5, and Step 8 shows selected/run typing survives a structured RG/measure toy by relocating the obstruction to the selection. R4 transparency was added in Step 5, and Step 6 shows the branch-split readouts are not removed by non-coordinate or nonlinear readouts of the endpoint source. R6 is the stress battery itself, Steps 6--8. The framework premise \(L=W\) restricted to regimes and the QM-GR interpretation documents are held.

Sources: \texttt{steps/step5\_honest\_regrade\_artifacts/results\_summary.md}, \texttt{steps/step6\_shared\_frame\_robustness\_artifacts/results\_summary.md}, \texttt{steps/step7\_e042\_earned\_continuation\_artifacts/results\_summary.md}, \texttt{steps/step8\_e021\_rg\_measure\_robustness\_artifacts/results\_summary.md}.

\end{document}
